% !TeX root = ../../../manual.tex
\chapter{TikZ Engineering Diagrams}
\label{ch:tikz-engineering}

\section{Flowcharts}

OmniLaTeX provides node styles for standard flowchart shapes via
\texttt{omnilatex-tikz-core.sty}:

\begin{longtblr}[
    caption = {Flowchart node styles},
    label = {tab:flowchart-styles},
]{colspec = {l l}, rowhead = 1}
    \toprule
    Style      & Shape \\
    \midrule
    startstop  & Rounded rectangle (terminal) \\
    io         & Parallelogram (input/output) \\
    process    & Rectangle (process step) \\
    decision   & Diamond (conditional branch) \\
    \bottomrule
\end{longtblr}

\begin{verbatim}
\begin{tikzpicture}[
    node distance=1.5cm,
    startstop/.style={...},
    process/.style={...},
    decision/.style={...},
]
    \node (start) [startstop] {Start};
    \node (proc1) [process, below=of start] {Read input};
    \node (dec1)  [decision, below=of proc1] {Valid?};
    \node (proc2) [process, below=of dec1]   {Process data};
    \node (stop)  [startstop, below=of proc2]{Stop};

    \draw[-{Stealth}] (start) -- (proc1);
    \draw[-{Stealth}] (proc1) -- (dec1);
    \draw[-{Stealth}] (dec1)  -- node[right] {Yes} (proc2);
    \draw[-{Stealth}] (dec1.west) -- ++(-2,0)
                      node[above] {No} |- (proc1);
    \draw[-{Stealth}] (proc2) -- (stop);
\end{tikzpicture}
\end{verbatim}

Arrow styles use the \texttt{arrows.meta} library. Common tips include
\texttt{Stealth}, \texttt{Latex}, \texttt{>}, and \texttt{|}.

\section{Circuit Diagrams}

The \texttt{circuits.ee.IEC} TikZ library provides standard electrical
symbols:

\begin{verbatim}
\usetikzlibrary{circuits.ee.IEC}

\begin{tikzpicture}[circuit ee IEC]
    \draw (0,0) to[resistor] (3,0)
               to[battery]   (3,-2)
               to[contact]   (0,-2) -- cycle;
\end{tikzpicture}
\end{verbatim}

\begin{longtblr}[
    caption = {IEC circuit symbols},
    label = {tab:iec-symbols},
]{colspec = {l l}, rowhead = 1}
    \toprule
    Symbol    & Component \\
    \midrule
    resistor  & Resistor \\
    capacitor & Capacitor \\
    inductor  & Inductor \\
    battery   & Voltage source / battery \\
    contact   & Switch / contact \\
    diode     & Diode \\
    \bottomrule
\end{longtblr}

Component values are specified with the \texttt{info} key:
\texttt{to[resistor, info=\$R\_1\{=\}100\,\textbackslash Omega\$]}.

\section{3D Drawings}

The \texttt{tdplot} library provides three-dimensional coordinate frames:

\begin{verbatim}
\usepackage{tikz-3dplot}
\tdplotsetmaincoords{70}{110}

\begin{tikzpicture}[tdplot_main_coords]
    \draw[->] (0,0,0) -- (4,0,0) node[right] {$x$};
    \draw[->] (0,0,0) -- (0,4,0) node[above] {$y$};
    \draw[->] (0,0,0) -- (0,0,4) node[below left] {$z$};

    \draw[fill=blue!20] (0,0,0) -- (3,0,0) -- (3,2,0)
                        -- (0,2,0) -- cycle;
\end{tikzpicture}
\end{verbatim}

Canvas planes (\texttt{xy}, \texttt{yz}, \texttt{xz}) help project 2D
annotations onto 3D drawings. Flow arrows use the \texttt{-{Stealth}}
arrow tip with the \texttt{line width} key for visibility in perspective
views.

\section{Custom Thermodynamic Shapes}

OmniLaTeX defines custom TikZ shapes for thermodynamic and process
engineering diagrams:

\begin{longtblr}[
    caption = {Custom thermodynamic shapes},
    label = {tab:thermo-shapes},
]{colspec = {l l l}, rowhead = 1}
    \toprule
    Shape                 & Anchors & Description \\
    \midrule
    valve                 & in, out & Control valve \\
    pump                  & in, out & Circulation pump \\
    compressor            & in, out & Compressor \\
    heat exchanger        & in, out & Heat exchanger \\
    simple heat exchanger & in, out & Simple HX \\
    radiator              & top, right, exit, lower entry, upper entry, ventilation & Radiator with multiple connections \\
    sensor                &         & Sensor \\
    levelindicator        &         & Level indicator \\
    boiler                &         & Boiler \\
    equalizing tank       &         & Equalizing tank \\
    \bottomrule
\end{longtblr}

Use these shapes in node definitions:

\begin{verbatim}
\node[valve] (v1) at (0,0) {};
\node[pump, right=2cm of v1] (p1) {};
\draw[-{Stealth}] (v1.out) -- (p1.in);
\end{verbatim}

Shapes with \texttt{in}/\texttt{out} anchors connect naturally to pipe
segments. The \texttt{radiator} shape provides multiple connection points
(\texttt{top}, \texttt{right}, \texttt{exit}, \texttt{lower entry},
\texttt{upper entry}, \texttt{ventilation}) for complex piping layouts.

\section{Pipe Networks}

Pipe networks use dedicated node styles for fittings and junctions:

\begin{verbatim}
\begin{tikzpicture}
    \node[pipe] (p1) at (0,0) {};
    \node[threeway valve] (tv1) at (3,0) {};
    \node[control valve]  (cv1) at (6,0) {};
    \node[fourway valve]  (fv1) at (3,-3) {};

    \draw[-{Stealth}] (p1) -- (tv1);
    \draw[-{Stealth}] (tv1) -- (cv1);
    \draw[-{Stealth}] (tv1) -- (fv1);
\end{tikzpicture}
\end{verbatim}

\begin{longtblr}[
    caption = {Pipe network node styles},
    label = {tab:pipe-styles},
]{colspec = {l l}, rowhead = 1}
    \toprule
    Style             & Description \\
    \midrule
    pipe              & Straight pipe segment \\
    threeway valve    & Three-way branching valve \\
    control valve     & Two-way control valve \\
    fourway valve     & Four-way cross valve \\
    \bottomrule
\end{longtblr}

Pipes connect via the \texttt{in}/\texttt{out} anchors on valve nodes.
The \texttt{pipe} node acts as a directional segment for flow
visualisation.

\section{Control System Diagrams}

Simulink-style block diagrams use rectangular nodes with named ports:

\begin{verbatim}
\begin{tikzpicture}[
    block/.style={draw, rectangle, minimum height=2.5em,
                  minimum width=4em, anchor=west},
    sum/.style={draw, circle, node distance=1.5cm},
]
    \node[sum]   (sum)  {$\Sigma$};
    \node[block, right=of sum]  (plant)  {Plant};
    \node[block, right=of plant](ctrl)   {Controller};
    \node[block, below=of ctrl] (sat)    {Saturation};

    \draw[-{Stealth}] (sum)  -- (plant);
    \draw[-{Stealth}] (plant) -- (ctrl);
    \draw[-{Stealth}] (ctrl)  -- (sat);
    \draw[-{Stealth}] (ctrl.east) -- ++(1,0) |- (sum.south)
        node[pos=0.95, left] {$-$};
\end{tikzpicture}
\end{verbatim}

Feedback loops are drawn by routing the output signal back to the summing
junction with \texttt{-{}- ++(dx,0) |-} paths. The saturation block models
actuator limits with \texttt{min}/\texttt{max} parameters.

\section{Annotation Arrows}

OmniLaTeX provides utility commands for engineering annotations:

\begin{longtblr}[
    caption = {Annotation arrow commands},
    label = {tab:annotation-arrows},
]{colspec = {l l}, rowhead = 1}
    \toprule
    Command                  & Description \\
    \midrule
    \texttt{\textbackslash annotationarrow}  & Arrow with a text label at the tip \\
    \texttt{\textbackslash wall}             & Hatched wall for boundary conditions \\
    \texttt{\textbackslash flowarrow}        & Directional flow arrow on a pipe or duct \\
    \texttt{\textbackslash arrowlabel}       & Text label positioned along an arrow \\
    \bottomrule
\end{longtblr}

\begin{verbatim}
\flowarrow{(0,0)}{(4,0)}{Flow direction}
\wall{(6,-1)}{(6,1)}
\annotationarrow{(2,2)}{(5,3)}{Note: pressure drop}
\end{verbatim}

These commands simplify common annotation patterns in process and
mechanical engineering diagrams, keeping the TikZ source concise and
readable.
